% ED26-SRC-027-501 separated continuation. Released core/marked-tree files frozen.
\subsection{The direct-sum theorem on genuinely mixed terms}
\label{sec:p12-klop1987-direct-sum-proof}

The 1987 tutorial's Theorem2.3.2 concerns ordinary first-order TRSs
whose \emph{whole nonvariable signatures} are disjoint, not merely
different rule names or different head constants
\parencite[pp.15--17]{Klop1987CWITutorial}. The following full
mixed-term proof reconstructs the preserved-layer and deletion-depth
mechanism of Toyama's separately identified author thesis
\parencite[\S2.2, pp.19--44]{Toyama1990DirectSumProof}.
This later proof source is not treated as a byte-identical copy of
the cited1987 publisher article or as the 1992 handbook chapter.

\begin{hbdefinition}{def:p12-klop1987-ds-signatures}{The exact unrestricted first-order regime}
Let \(\Sigma_1\cap\Sigma_2=\varnothing\), including constants.
Each operator has a fixed finite arity. Variables belong to one
common infinite supply \(V\). Each rule of \(R_i\) has a
nonvariable left side over \(\Sigma_i\), and its right-side
variables occur on its left. Terms are finite free terms, and
rewriting is closed under arbitrary first-order substitutions and
contexts. No left-linearity, nonduplication, noncollapsing condition
or termination assumption is imposed. The direct sum is
\(R=R_1\oplus R_2\) on \(T(\Sigma_1\cup\Sigma_2,V)\).
If signatures initially share a name, take renamed isomorphic
copies before forming this sum; an ordinary shared-signature union
is a different operation.
For the forward proof assume CR of both components on \emph{all}
their terms, not just ground terms.
\end{hbdefinition}

\begin{hbdefinition}{def:p12-klop1987-ds-layers}{Caps, principal aliens, rank and preservation}
The colour of a nonvariable root is its signature index. At a root
of colour \(a\), the maximal connected \(\Sigma_a\)-cap is
obtained by cutting every first opposite-colour subterm on each
path. These cut subterms are its principal aliens. Variables are
neutral leaves. Write \(t=C[A_1,\ldots,A_k]\) for this occurrence
decomposition; repeated values still have separate occurrence holes.
Its layer rank is1 for a pure one-signature term or variable, and
\(1+\max_i\operatorname{rank}(A_i)\) for a mixed term.
The special subterms are the root term and the recursively obtained
principal aliens.

Pure one-signature terms and variables are called preserved.
A mixed term of root colour \(a\) is preserved when every reduct
retains a root in \(\Sigma_a\), and all its principal aliens are
preserved. The universal root condition is a mathematical predicate,
not an effective normalizing oracle. It is used below only for an
existence proof; terms failing it are handled by a separate finite
deletion construction.

For the cap proof also use a frontier presentation \(t=H[\mathbf A]\)
that replaces variable leaves by holes as well as principal aliens.
Thus \(H\) contains only operators of one fixed signature and
holes. An empty cap may be a single hole. This presentation is
explicitly different from the principal-alien tuple used for layer rank.
\end{hbdefinition}

\begin{hblemma}{thm:p12-klop1987-ds-invariants}{proved=annex:p12-klop1987-ds-invariants}{Layer invariants without termination}
Rewriting cannot increase layer rank. Preserved terms remain
preserved under reduction. A preserved alien whose initial root
is outside \(\Sigma_a\) never acquires a \(\Sigma_a\)-root;
a pure alien may become a neutral variable, which also remains
outside that signature. These assertions do not bound reduction
length or term size.
\end{hblemma}

\begin{hbdefinition}{def:p12-klop1987-ds-coherence}{The equality condition required by nonlinear patterns}
For equal-length finite tuples write \(\mathbf A\propto\mathbf B\)
if \(A_i=A_j\Rightarrow B_i=B_j\) for all indices. Equality
is literal finite-term equality. Tuple reduction is componentwise.
Coherence does not require distinct inputs to remain distinct.
For a term \(u\), CR at \(u\) means that any two finite
reductions from it join; it is inherited by its descendants.
\end{hbdefinition}

\begin{hblemma}{thm:p12-klop1987-ds-coherence}{proved=annex:p12-klop1987-ds-coherence}{Finite common frontiers and exact cap-step transport}
If each member of a finite frontier tuple has CR, finitely many
chosen descendant tuples have a further common descendant tuple
coherent with each of them and with the original equalities.
A monochromatic cap rewrite whose pattern cannot enter the frontier
values transports to a coherently updated tuple, with the same
rule and the same copied/erased index occurrences on its right.
This covers nonlinear left sides, duplicating right sides and
collapsing right sides.
\end{hblemma}

\begin{hbexample}{def:p12-klop1987-ds-not-commutation}{Why the original component relations need not commute}
Take \(R_1=\{D(x,x)\to E\}\) and \(R_2=\{a\to b\}\),
with disjoint signatures. From \(D(a,a)\), the first component
gives \(E\), while the second gives \(D(b,a)\). No pure
\(R_1\)-path from \(D(b,a)\) joins a pure \(R_2\)-path
from \(E\): neither endpoint admits the needed component step.
The sum joins by \(D(b,a)\to_2D(b,b)\to_1E\).
The tuple update \((a,a)\to(b,a)\) was not coherent; its
further update to \((b,b)\) is. The proof therefore repairs
frontiers, rather than asserting commutation of arbitrary unions.
\end{hbexample}

\begin{hblemma}{thm:p12-klop1987-ds-preserved}{proved=annex:p12-klop1987-ds-preserved}{Preserved caps and their collapsing boundary}
A term consisting of a fixed-colour cap whose frontier aliens are
preserved and have CR itself has CR. Its outer cap may collapse
to one alien. Consequently every preserved term has CR, by
induction on layer rank. The cap statement also proves CR of a
nonpreserved term all of whose principal aliens are preserved.
\end{hblemma}

\begin{hbdefinition}{def:p12-klop1987-ds-depth}{Minimal nonpreserved layers and parallel deletion}
A deletion occurrence is a nonpreserved special subterm all of
whose principal aliens are preserved. Choose a finite reduction
of it to the first term whose root is outside its initial signature.
This root-changing reduct is preserved. A deletion replaces the
occurrence by that reduct, using an actual finite \(R\)-path.
Deletion occurrences are disjoint, because none lies below a
preserved principal alien of another.

Let \(d(t)=0\) for a preserved term; \(d(t)=1\) when \(t\)
itself is a deletion occurrence; otherwise, for a nonpreserved
mixed \(t=C[A_1,\ldots,A_k]\), set
\(d(t)=1+\max_i d(A_i)\). This is the maximum layer depth
of its minimal nonpreserved occurrences. A parallel deletion
\(t\Rightarrow_{
m d}u\) deletes all occurrences at that
maximum depth, with independently chosen finite witnesses.
No computable selector for this semantic classification is claimed.
\end{hbdefinition}

\begin{hblemma}{thm:p12-klop1987-ds-depth}{proved=annex:p12-klop1987-ds-depth}{The actual decreasing depth}
The depth is a finite natural number. Ordinary rewriting cannot
increase it. Every nonpreserved term has a parallel deletion;
such a deletion is a finite ordinary reduction and strictly
decreases \(d\). Every deletion occurrence has CR by the cap
lemma. This does not assert termination of ordinary rewriting.
\end{hblemma}

\begin{hblemma}{thm:p12-klop1987-ds-diagrams}{proved=annex:p12-klop1987-ds-diagrams}{All ordinary-step/parallel-deletion positions}
Given \(t\to u\) and \(t\Rightarrow_{
m d}v\), with
\(d(t)>0\), the following joins exist.
If \(d(u)=d(t)\), one can choose
\(u\Rightarrow_{
m d}w\), with \(v\to^*z,w\to^*z\)
and \(d(z)<d(t)\). If \(d(u)<d(t)\), one can choose
\(v\to^*z,u\to^*z\) with that same strict bound.
The proof treats disjoint steps, a step inside a selected layer,
and a step above selected layers. In the last case coherent
joining is essential when copies reduce differently.
\end{hblemma}

\begin{hbtheorem}{thm:p12-klop1987-ds-confluence}{proved=annex:p12-klop1987-ds-confluence}{The tutorial's unrestricted direct-sum theorem}
For the exact regime above,
\[
 \mathrm{CR}(R_1\oplus R_2)
 \quad\Longleftrightarrow\quad
 \mathrm{CR}(R_1)\ \mathbin{\&}\ \mathrm{CR}(R_2).
\]
This includes mixed terms, nonlinear patterns, duplication and
collapse. It does not assume either component terminates, and
does not replace the sum by a shared-signature union.
\end{hbtheorem}

\begin{hbtheorem}{thm:p12-klop1987-ds-first-counter}{proved=annex:p12-klop1987-ds-first-counter}{The source's first termination failure}
Let \(R_1=\{f(0,1,x)\to f(x,x,x)\}\) and
\(R_2=\{\mathsf{or}(x,y)\to x,\mathsf{or}(x,y)\to y\}\),
over disjoint signatures. Both terminate. With
\(q=\mathsf{or}(0,1)\), their sum has the cycle
\[
 f(q,q,q)\to f(0,q,q)\to f(0,1,q)\to f(q,q,q).
\]
Here \(R_2\) is not CR. This is the particular source counterexample,
not an argument that two complete factors can never be combined.
\end{hbtheorem}

\begin{hbdefinition}{def:p12-klop1987-ds-complete-counter}{The exact eleven-rule/four-argument source factor}
In the next source example the glyphs1 through7 are distinct
nullary symbols, not numbers in a semantic ordering. The first
factor has the two rules
\(F(4,5,6,x)\to F(x,x,x,x)\),
\(F(x,y,z,w)\to7\), and all nine constant rules
\[
1\to4,\ 1\to5,\ 2\to4,\ 2\to6,\ 3\to5,\ 3\to6,
\quad4\to7,\ 5\to7,\ 6\to7.
\]
The disjoint second factor has exactly
\(G(x,x,y)\to x\), \(G(x,y,x)\to x\),
\(G(y,x,x)\to x\).
These are the whole source graph and its three majority rules.
\end{hbdefinition}

\begin{hbtheorem}{thm:p12-klop1987-ds-complete-counter}{proved=annex:p12-klop1987-ds-complete-counter}{Completeness is also not modular without extra hypotheses}
Both factors in the preceding definition are SN and CR. Their
sum is CR by the full modularity theorem, but is not SN.
Putting \(q=G(1,2,3)\), the source cycle is
\begin{align*}
 F(q,q,q,q)&\to^*
 F(G(4,4,3),G(5,2,5),G(1,6,6),q)\\
 &\to^*F(4,5,6,q)\to F(q,q,q,q).
\end{align*}
The source attributes this construction to Barendregt and Klop,
adapted by Toyama. It refutes modularity of SN and of CR plus SN;
it does not refute confluence modularity.
\end{hbtheorem}

\begin{hbtheorem}{thm:p12-klop1987-ds-applications}{proved=annex:p12-klop1987-ds-applications}{The binary equality test and selector applications}
The one-rule binary factor \(D(x,x)\to E\) is SN and CR.
The factor with \(\mathsf{if}(\mathsf{true},x,y)\to x\),
\(\mathsf{if}(\mathsf{false},x,y)\to y\),
\(\mathsf{if}(z,x,x)\to x\) is also SN and CR.
Their separate direct sums with CL are therefore CR. The nonlinear
test and selector do not become regular by being disjoint factors.
Here \(\mathsf{if}\) is an ordinary ternary symbol, not a condition
attached to a rule.
\end{hbtheorem}

The report separately contrasts the applicative test \(Dxx\to E\).
Its restored left side is
\(\operatorname{Ap}(\operatorname{Ap}(D,x),x)\), so that
extension shares the application operator with CL. It is not this
direct sum; the source reports it nonconfluent and refers its proof
to Klop(1980). That separate nonconfluence witness is not supplied
by the binary-factor proof or by a counterexample to an unrelated union.
The report also mentions reduced/irreducible variants in Toyama's
1986 preprint and Ganzinger--Giegerich(1987), and reports the latter's
extra \(b\to k(0)\) arrow as producing a self-reduction of
\(f(b,b,b,g(0,b))\). These cited variants retain their own rule
sets and proof standing, rather than being added to the eleven-rule
factor above.
